\documentclass[11pt]{article}
\usepackage{amsmath,amssymb}
\title{A leading kinetic bound for nonconstant harmonic shear fields}\author{}\date{}
\begin{document}\maketitle
\paragraph{An actual harmonic family beyond constant curvature.}
On the original cube of side $R\ge2$, put
\[
 a_1(x)=A(x_0,x_2,x_3)\,i\sigma_3,\qquad a_\mu=0\quad(\mu\ne1),
                                                               \tag{1}
\]
where $A$ is discrete harmonic on the interior of the three-dimensional
transverse cube, with its complete boundary values retained.
The original linearized Maxwell force and divergence of (1) vanish
on all free links. It is the true harmonic extension of its original
boundary tangent. Its energy is
\[
 \mathcal E(a)=\sum_p |(d_1a)_p|^2
       =R\sum_{\text{transverse edges}}|\delta A|^2.              \tag{2}
\]
Fixed-only plaquettes are included.

For sufficiently small physical source parameter $\epsilon$, the actual
nearby tree-quotient action minimum exists uniquely at each fixed $R$.
Cartan conjugation and sign reversal imply an odd Cartan minimum branch.
Its first derivative is (1) and its second derivative is zero. Thus
the leading determinant's second derivative is computed from the
linear background $\epsilon a$ without missing an acceleration term,
even though this background need not be an exact nonlinear minimum
at nonzero $\epsilon$.

\paragraph{1. Literal scalar connections from the full Wilson Hessian.}
Use the full charged vector operator $V=C+DD^*$ and ghost $Q=D^*D$
of the constant flux spin and magnetic reduction calculation, before tree deletion. On a component grid $\mu\ne1$,
the scalar same-component connection is
\[
 a^{(\mu)}_1(x)=A(x)+A(x+\hat\mu),\qquad
 a^{(\mu)}_\nu=0\quad(\nu\ne1).
                                                               \tag{3}
\]
On component one its raw transverse phases are
$\delta_\nu A$ and its direction-one phase is $2A$.
The exact scalar gradient gauge $A(x)$ removes those transverse phases,
leaving $a^{(1)}_1=2A$. The ghost likewise has $a^{(g)}_1=2A$.
These follow from the original positive-after/inverse-before
insertions and endpoint gauge factors, not a continuum substitution.
All these connections are independent of coordinate one.

The component-one grid has Neumann direction one and three Dirichlet
transverse directions. Every OTHER component and the ghost has
Dirichlet direction one. This precise boundary fact is used below.
The transverse nearest-neighbor energy of every coefficient in (3),
times its number of direction-one sites, is at most $C\mathcal E(a)$
by a literal bounded-multiplicity sum of original plaquette differences.

\paragraph{2. Sum open-current moments before closing the loop.}
First consider a scalar grid with DIRICHLET coordinate $Y=X^1$ and
connection $b(Z)dY$, independent of $Y$.
For each start $(y,z)$ use the constant-gradient gauge subtracting $b(z)$.
The current
\[
 J_t=\int_0^t[b(Z_{s-})-b(z)]\,dY_s
\]
stopped INCLUDING its possible killing jump is a martingale.
For the transverse coordinate process, domination by the reflected
chain on the same finite set gives
\[
 \sum_z E_z[(b(Z_s)-b(z))^2;\text{survival}]
 \le2\langle b,(I-e^{-sL_N})b\rangle
 \le2s\,\mathcal E_N(b).
\]
Here $L_N$ is used ONLY for this moment upper bound, with its original
internal edge energy and suppressed outward jumps; the actual path
kernel still has its killed boundary conditions.
Summing starts and integrating the quadratic variation yields
\[
 \sum_x E_x[J_t^2;\tau>t]\le C t^2\mathcal E(a).                  \tag{4}
\]
No pointwise bound on $A$ or its boundary values is needed.

For component one, $Y$ reflects and all three coordinates $Z$ are
Dirichlet. Its coefficient $b=2A$ is harmonic. Use exact discrete
integration by parts to write the SAME centered-gauge current as
\[
 J_t=-\int_0^t(Y_{s-}-y)\,db(Z_s)
                   +(Y_t-y)(b(Z_t)-b(z)).                      \tag{5}
\]
The first term is a stopped martingale: its drift is
$-(Y-y)\Delta b$, which vanishes at every interior transverse site.
For a killing jump use the ACTUAL boundary value of $b$; this retains
the harmonic equation, and does not change any surviving path current.
Its summed quadratic variation is bounded by
$Ct^2 n_Y\sum|\delta b|^2$.
For the endpoint product in (5), independence of the reflected $Y$
and transverse killed processes, the displacement bound
$E(Y_t-y)^2\le2t$, and the preceding reflected-energy estimate give
the same upper bound. Thus (4) holds also for this component.
The killing and endpoint terms have not been discarded.

\paragraph{3. Scalar determinant curvature.}
The original scalar kernels satisfy
$\sup p_t\le C(1+t)^{-2}e^{-ct/R^2}$.
The SAME-start half-loop reversal of the mixed boundary magnetic loop bound calculation, now using the summed
moment (4), proves
\[
 -\partial_\epsilon^2\operatorname{Tr}e^{-tL_\epsilon}\big|_0
 \le C\,t^2(1+t)^{-2}e^{-ct/R^2}\mathcal E(a).
\]
Integrating $dt/t$ gives, for EACH original scalar grid,
\[
 0\le(\log\det L_\epsilon)''_0
                  \le C\log(2+R)\,\mathcal E(a).                \tag{6}
\]
The complete boundary conditions and all scalar gauge transforms
are retained. Harmonicity in (5) is essential; an arbitrary field
with a reflected current coordinate is not covered by this proof.

\paragraph{4. Spin and diagonal partners with nonconstant curvature.}
Let $f_p=(d_1a)_p$. The exact charged decomposition is still
$V=L+P+S$: $L$ is component diagonal, $P$ has diagonal
$\sum_{p\ni e}(\cos(\epsilon f_p)-1)$, and $S$ is off component.
At zero, $S'$ is a bounded-stencil matrix whose entries are signed
local fluxes $f_p$. The mixed trace
$\operatorname{Tr}(GL'GS')$ and $\operatorname{Tr}(GS'')$ vanish
by component structure exactly as in the constant flux spin and magnetic reduction calculation.

The original flat product Green kernels obey
\[
 |G_\mu(x,y)|\le C(1+|x-y|)^{-2}.                               \tag{7}
\]
One direct finite-boundary proof uses the one-dimensional reflected
image formula for an upper bound at $t\le R^2$, while Dirichlet killing
can only decrease the positive kernel. The unrestricted kernel satisfies
\[
 p_t^{\mathbb Z}(n)\le C(1+t)^{-1/2}
  e^{-c\min(n^2/(1+t),|n|)}
\]
by shifting its elementary Fourier integral by
$\min(1,|n|/(4(1+t)))$ in the imaginary direction.
The reflected image sums have periods comparable to $R$ and uniformly
bounded multiplicity at $t\le R^2$. Their four-dimensional product gives
the corresponding off-diagonal bound. For $t\ge R^2$ the original
three-Dirichlet product gap gives $CR^{-4}e^{-ct/R^2}$.
Integrating these two ranges proves (7), including boundary rows.

Consequently the absolute spin trace is bounded by a kernel
$C(1+|p-q|)^{-4}$ on local plaquette fluxes:
\[
 |\operatorname{Tr}(GS'GS')|
 \le C\sum_{p,q}\frac{|f_p f_q|}{(1+|p-q|)^4}
 \le C\log(2+R)\sum_p f_p^2.                                   \tag{8}
\]
The last step is the finite four-dimensional shell sum and Schur
bound, not a commutativity assumption.
The diagonal correction costs at most $C\sum f_p^2$ by the bounded
Green diagonal. The neutral term is
$-\frac12\sum_p f_p^2 r_p G r_p^*$; each $r_p G r_p^*\le1$
because $r_p^*r_p\le H\le\Delta_1$.
All full gauge/ghost determinant multiplicities remain those of the constant flux spin and magnetic reduction calculation.

Combining these ACTUAL scalar, spin, diagonal and neutral partners yields
\[
 \boxed{|F_{\rm quad}''[a,a]|
       \le C\log(2+R)\mathcal E(a),\qquad
       |h_a^*\widehat K h_a|
       \le C\log(2+R)\mathcal E(a).}                            \tag{9}
\]
The second identity uses the evaluated harmonic matching of the complete harmonic kinetic matching calculation.
In canonical source units multiply (9) by $B^{-1}$.
This is a leading, complete, UNCUT kinetic bound on the harmonic family
(1), not a two-loop or original nonlinear remainder bound.

\paragraph{A mode completely outside the eighteen constant-curvature modes.}
Take
\[
 A(x_0,x_2,x_3)=
 \frac{(x_0-R/2)(x_2-R/2)}{R^3}.
\]
It is exactly discrete harmonic, its mean flux in every plane is zero,
and
\[
 \mathcal E(a)=\frac{(R+1)^2(R+2)}{6R^3}\longrightarrow\frac16.
\]
Thus (9) controls a genuine nonconstant harmonic boundary mode with
ZERO projection onto the eighteen modes of the complete constant curvature coarse sector calculation.
The exact original-incidence the finite calculation verifies its harmonic
force, zero coarse flux, energy and literal scalar phases on $R=2,3$.

Flat global color invariance also gives a precise corollary: within
ONE spatial polarization, $a_1=\sum_a A^a(Z)i\sigma_a$ with each $A^a$
harmonic obeys (9) with the sum of their energies. The Hessian is
diagonal in color. This argument concerns the second derivative at
flat background and does not assume that its finite noncommuting
source is a Cartan configuration.

The family (1) grows with the number of transverse boundary values.
It is not the entire harmonic boundary space: superposing different
polarizations or arbitrary multi-component sources requires further mixed
estimates. The global self-energy is not reset to independent subblock
tensors. The transverse two-loop packet, its ghost/cutoff/source partners
and the original uniform remainder remain open.
\end{document}
